\documentclass[10pt]{article}
\usepackage[letterpaper,margin=0.68in]{geometry}
\usepackage{amsmath,amssymb}
\usepackage{enumitem}
\usepackage{xcolor}
\usepackage{fancyhdr}
\usepackage{microtype}
\usepackage[most]{tcolorbox}
\definecolor{olive}{HTML}{4D5430}
\definecolor{gold}{HTML}{9A7B22}
\definecolor{pale}{HTML}{F5F1DC}
\definecolor{softgreen}{HTML}{EDF0E2}
\pagestyle{fancy}\fancyhf{}\lhead{\textcolor{olive}{MATH\& 107 Answer Key}}\rhead{\textcolor{gold}{Fall 2026}}\cfoot{\thepage}
\setlength{\headheight}{14pt}
\setlength{\parindent}{0pt}\setlength{\parskip}{5pt}
\setlist[enumerate]{leftmargin=*,itemsep=5pt,topsep=4pt}
\newtcolorbox{solbox}[1]{colback=pale,colframe=olive,boxrule=.7pt,arc=2mm,title=\textbf{#1},fonttitle=\color{olive},left=2mm,right=2mm,top=1.5mm,bottom=1.5mm}
\newcommand{\modulekey}[2]{\clearpage{\Large\bfseries\textcolor{olive}{Module #1: #2 — Answer Key}}\par\textcolor{gold}{\rule{\linewidth}{1.2pt}}\par}
\begin{document}
\modulekey{2}{Frequency, Probability, and Distributions}
\textbf{Prequiz}
\begin{enumerate}
\item Population: \textbf{all students at the college}. Sample: \textbf{the 150 students in classes before 9 a.m.} Students who do not take classes before 9 a.m. are not represented by this selection rule. If early-class students differ systematically from other students in work schedules, the sample may not represent the whole college well.
\item Bus:
\[
\frac6{24}=\boxed{0.25}.
\]
Not driving means bus, bike, or walk:
\[
\frac{6+4+2}{24}=\frac{12}{24}=\boxed{0.50}.
\]
A \textbf{bar chart} is appropriate because transportation method is categorical.
\item A \textbf{histogram} is appropriate because the ages are quantitative values grouped into numerical intervals.
\item Known probabilities sum to
\[
0.18+0.27+0.31=0.76,
\]
so
\[
p=1-0.76=\boxed{0.24}.
\]
The complement of the \(0.18\) event is
\[
1-0.18=\boxed{0.82}.
\]
\item
\[
P(\text{sum}=7)=\frac6{36}=\frac16\approx\boxed{0.1667}.
\]
\end{enumerate}

\textbf{Matching quiz}
\begin{enumerate}
\item Population: all registered county voters. Sample: the \(500\) selected voters. The intended strength of the plan is \textbf{random selection from the population of interest}, which is designed to give a broader cross-section than a convenience sample from one time or place.
\item
\[
\frac{18}{60}=\boxed{0.30}.
\]
\item \textbf{Histogram}, because commute time is numerical and grouped into intervals.
\item
\[
p=1-(0.22+0.35+0.19)=1-0.76=\boxed{0.24}.
\]
Complement of the \(0.35\) event:
\[
1-0.35=\boxed{0.65}.
\]
\item
\[
P(\text{blue})=\frac38=\boxed{0.375},\qquad P(\text{not blue})=\frac58=\boxed{0.625}.
\]
\end{enumerate}
\end{document}
